\documentclass[10pt,a4paper]{article}
\usepackage[utf8]{inputenc}
\usepackage{amsmath}
\usepackage{amsfonts}
\usepackage{amssymb}
\usepackage{xcolor, soul}
\sethlcolor{green}
\usepackage{enumerate}
\title{GSF-3100 \\
02-Mathématiques financières des obligations \\
Série d'exercices 1 (Solutions)}
\begin{document}
\maketitle

\newpage

\section*{Exercice 1}
Trouver le taux effectif semestriel équivalent à un taux nominal annuel de 9 \% capitalisé en continu.

\begin{enumerate}[a)]
\item 4,65\%
\item 4,55\%
\item \hl{4,60\%}
\item 4,70\%
\end{enumerate}

\section*{Solution 1}

\textbf{Paramètres:}
\begin{itemize}
\item $u = 1$ période de 1 semestre
\item $v = 2$ périodes de 1 semestre par année
\end{itemize}

\textbf{Taux effectif annuel:}
\begin{align*}
e^{0.09}-1 = 9.4174 \%
\end{align*}

\textbf{Taux effectif semestriel:}
\begin{align*}
r_1 = (1+r_2)^{\frac{u}{v}}-1
\end{align*}

\begin{align*}
r_{semestriel} = (1.094174)^{\frac{1}{2}}-1 = 4.6028 \%
\end{align*}

\newpage


\section*{Exercice 2}
Trouver le taux nominal quotidien capitalisé annuellement équivalent à un taux nominal bisannuel de 20 \% capitalisé trimestriellement.

\begin{enumerate}[a)]
\item 0,045\%
\item 0,015\%
\item 0,022\%
\item \hl{0,028\%}
\end{enumerate}

\section*{Solution 2}


Taux effectif trimestriel équivalent au taux nominal connu ($\frac{1}{8}$ de deux ans avant d'effectuer la première capitalisation):
\begin{align*}
r_{trimestriel} = 0.20 \times \frac{1}{8} = 2.5 \% \textbf{ par trimestre}
\end{align*}


Taux effectif trimestriel (connu) $\rightarrow$ Taux effectif annuel (recherché)\\

\textbf{Paramètres:}
\begin{itemize}
\item $r_2 = 0,025$ 
\item $u = 1$ période de 1 an
\item $v = \frac{1}{4} = 0,25$ période de 1 an par trimestre.
\end{itemize}

\begin{align*}
r_1 = (1+r_2)^{\frac{u}{v}}-1
\end{align*}

\begin{align*}
r_1 = (1+0.025)^{\frac{1}{\frac{1}{4}}}-1
\end{align*}

\begin{align*}
r_1 = (1+0.025)^4-1
\end{align*}

\begin{align*}
r_1 =r_{annuel} (1+0.025)^4-1=10.381\%
\end{align*}
\textbf{Note:} \\ On veut un taux effectif annuel car le taux nominal recherché est capitalisé annuellement.

\newpage

Taux effectif annuel $\rightarrow$ Taux nominal quotidien capitalisé annuellement

\begin{align*}
i=\frac{r_{annuel}}{m}
\end{align*}

\begin{align*}
i =\frac{r_{annuel}}{m} = \frac{0.10381}{365} = 0,028\% 
\end{align*}
\textbf{Note:} \\ $m$ correspond à la durée (quotidienne) du taux nominal recherché exprimé en fréquence de la période de capitalisation, soit (1/365) année pour une durée de 1 jour.

\newpage

\section*{Exercice 3}
Trouver la valeur actuelle d’une annuité de 35 paiements hebdomadaires de 3 \$ si le taux nominal annuel capitalisé hebdomadairement est de 12 \%.

\begin{enumerate}[a)]
\item 100.35\$
\item 100.79\$
\item 101.40\$
\item \hl{100.75\$}
\end{enumerate}

\section*{Solution 3}

Taux nominal annuel cap. hebdomadairement $\rightarrow$ Taux effectif hebdomadaire

\begin{align*}
r_{hebdo} = \frac{0.12}{52} = 0.00231
\end{align*}

Valeur actuelle d'une annuité:\\

\textbf{Paramètres:}
\begin{itemize} 
\item $A = 3$ \$ 
\item $r_{hebdo} = 0.00231$
\item $n = 35$ semaines
\end{itemize} 

\begin{align*}
VP_0=A \times \left[\frac{1-\left( \frac{1}{(1+r)^n}\right)}{r} \right] 
\end{align*}

\begin{align*}
VP_0=3 \times \left[\frac{1-\left( \frac{1}{(1+0,00231)^{35}}\right)}{0,00231} \right] 
\end{align*}

\begin{align*}
VP_0=3 \times 33,58527
\end{align*}


\begin{align*}
VP_0 = 100,75581 \$
\end{align*}


\newpage



\section*{Exercice 4}
Trouver la valeur future d’une annuité de 2 paiements semestriels de 5 \$ en supposant que le taux nominal semestriel capitalisé trimestriellement est de 7 \% et que le premier paiement aura lieu aujourd'hui (début de période).

\begin{enumerate}[a)]
\item \hl{11.09\$}
\item 10.99\$
\item 12.09\$
\item 11.40\$
\end{enumerate}

\section*{Solution 4}

Taux nominal semestriel capitalisé trimestriellement $\rightarrow$ Taux effectif trimestriel

\begin{align*}
r_{trimestriel} = \frac{0.07}{2} = 0.035
\end{align*}

Taux effectif trimestriel $r_2$ $\rightarrow$ Taux effectif semestriel $r_1$ (qui correspond à la fréquence des paiements)

\begin{align*}
r_1 = (1+r_2)^{\frac{u}{v}} -1
\end{align*}

\begin{align*}
r_1 = (1+0.035)^{\frac{2}{1}} -1
\end{align*}

\begin{align*}
r_1 = 0,07123
\end{align*}

Valeur future d'une annuité de début de période\\

Paramètres:
\begin{itemize}
\item $A = 5$ \$
\item $n = 2$ semestres
\item $r_{semestriel}=0,07123$
\item $k=1$ semestre "en avance"
\end{itemize} 


\begin{align*}
VF_{n+k} = A \times \left[ \frac{(1+r)^n-1}{r} \right] \times (1+r)
\end{align*}

\begin{align*}
VF_{2+1} = 5 \times \left[ \frac{(1+0,07123)^2-1}{0,07123} \right] \times (1+0,07123)
\end{align*}

\begin{align*}
VF_{3} = 5 \times 2.07123 \times (1+0,07123)
\end{align*}

\begin{align*}
VF_{3} =  11,09382 \$
\end{align*}
\textbf{Note:} \\ Puisque le premier paiement s'effectue aujourd'hui au lieu de la fin du premier semestre, tous les paiements suivants sont également anticipés d'une période à l'avance (k = 1) et rapportent des intérêts supplémentaires. Mathématiquement, c'est le terme (1+r)k qui permet d'incorporer cette hausse de valeur de l'annuité.

\newpage


\section*{Exercice 5}
Trouver le prix d’une obligation d’une valeur nominale de 1000 \$ venant à échéance dans 21,5 ans et ayant un taux de coupon (TC) de 24 \% annuel payable quatre fois par année, soit à chaque trimestre, en supposant un taux nominal annuel capitalisé mensuellement de 12 \%.

\begin{enumerate}[a)]
\item 1000\$
\item 1495.35\$
\item 1910.59\$
\item \hl{1904.9\$}
\end{enumerate}


\section*{Solution 5}
\begin{align*}
r_{mensuel}  = \frac{0.12}{12} = 0,01
\end{align*}

\begin{align*}
r_{trimestriel} = (1,01)^3-1 = 0,03030
\end{align*}

Paramètres:  
\begin{itemize}
\item $M = 1000$ \$
\item $n = 21,5$ ans $\times$ $4$ trimestres/an $= 86$ trimestres 
\item $r = 0,03030$
\item $C = \frac{TC}{4} \times M  = \frac{0.24}{4} \times 1000 = 60$ \$
\end{itemize}

\begin{align*}
P_0=C \times \left[\frac{1-\left( \frac{1}{(1+r)^n}\right)}{r} \right] +\frac{M}{(1+r)^n}
\end{align*}

\begin{align*}
P_0=60 \times \left[\frac{1-\left( \frac{1}{(1+0,03030)^{86}}\right)}{0,03030} \right] +\frac{1000}{(1+0,03030)^{86}}
\end{align*}

\begin{align*}
P_0 = 1828,2024 + 76,75769
\end{align*}

\begin{align*}
P_0 = 1904,96009
\end{align*}


\newpage

\section*{Exercice 6}
Trouver le prix d’une obligation perpétuelle (à vie) ayant un taux de coupon de 16 \% annuel payable deux fois par année et une valeur nominale de 1000 \$ en supposant un taux eectif semestriel de 6 \%.

\begin{enumerate}[a)]
\item 1433.3\$
\item \hl{1333.3\$}
\item 1336.3\$
\item 1400.3\$
\end{enumerate}


\section*{Solution 6}
Paramètres:
\item $M = 1000$ \$
\item $C = \frac{TC}{2} \times M = 0,08 \times 1000 = 80$ \$ 
\item $r = 6$ \%
\item $n = \infty$
\end{itemize}
\vspace{0.5cm}
\textbf{Note:} \\ Puisque le coupon est semestriel, nous utilisons le taux effectif semestriel tel que fourni dans l'énoncé.

\vspace{0.5cm}
Prix d'une obligation perpétuelle:
\begin{align*}
P_0 = \frac{C}{r}
\end{align*}

\begin{align*}
P_0 = \frac{80}{0.06}
\end{align*}

\begin{align*}
P_0 = 1333,33333 \$
\end{align*}
\textbf{Interprétation:} \\ Le prix d'une obligation qui verse un coupon de 80 \$ chaque semestre jusqu'à l'infini vaut 1333,33 \$ pour un investisseur qui exige un taux de rendement effectif semestriel de 6 \%. Cela pourrait par exemple s'agir d'une rente viagère ou de produits financiers liés à la planification d'une retraite.

\newpage

\section*{Exercice 7}
Trouver le prix $P_k$ d’une obligation d’une valeur nominale de 800 \$ venant à échéance dans 5,5 ans et ayant un taux de coupon de 10 \% annuel payable deux fois par année en supposant que le taux nominal annuel capitalisé semestriellement est de 12 \%.  Cette obligation se transigeait initialement à $P_0$ = 736,9050 \$ et le dernier coupon a été versé le 31 décembre dernier il y a exactement 126 jours. (Utiliser la méthode exact/exact.)

\begin{enumerate}[a)]
\item 787.1\$
\item 763.1\$
\item \hl{767.1\$}
\item 796.9\$
\end{enumerate}

\section*{Solution 7}
Taux effectif semestriel

\begin{align*}
r_{semestriel} = \frac{0.12}{2}= 0.06
\end{align*}

Valeur de k avec la méthode \textbf{exact/exact}
\begin{itemize}

\item  Nombre de jours exact qui s'est écoulé depuis le dernier coupon (numérateur): 126 jours

\item Nombre de jours exact de la période de coupon entre le 1er janvier et le 30 juin (dénominateur):
\begin{align*}
31+30+31+30+31+30 = 31 \times 3 + 30 \times 3 = 183 
\end{align*}

\item $k = \frac{126}{183} = 0,68852$ an
\end{itemize}

Prix de l'obligation à la période k \\

\textbf{Paramètres:}
\begin{itemize}
\item $M = 800$ \$
\item $n = 5,5$ ans $\times 2$ semestres/an $= 11$ semestres
\item $C = \frac{0.10}{2} \times 800 = 40$ \$
\item $r = 0,06$
\end{itemize}

\begin{align*}
P_k=P_0 \times (1+r)^k
\end{align*}

\begin{align*}
P_k= \left( C \times \left[\frac{1-\left( \frac{1}{(1+r)^n}\right)}{r} \right] +\frac{M}{(1+r)^n} \right) \times (1+r)^k
\end{align*}

\begin{align*}
P_k= \left( 40 \times \left[\frac{1-\left( \frac{1}{(1+0.06)^{11}}\right)}{0.06} \right] +\frac{800}{(1+0.06)^{11}} \right) \times (1+0.06)^{\frac{126}{183}}
\end{align*}

\begin{align*}
P_k = 767,0704 \$
\end{align*}

\newpage

\section*{Exercice 8}
Trouver la cote d'une obligation ayant un prix $P_k$ = 85 \$ entre deux dates de coupon et une fraction de période $k = 0,87$ écoulée depuis le versement du dernier coupon d'une valeur $C = 4$ \$.

\begin{enumerate}[a)]
\item 81.00\$
\item 81.13\$
\item 84.00\$
\item \hl{81.52\$}
\end{enumerate}

\section*{Solution 8}

Intérêt couru sur l'obligation
\begin{align*}
IC_k = k \times C
\end{align*}

\begin{align*}
IC_{0,87} = 0,87 \times 4 \$
\end{align*}

\begin{align*}
IC_{0,87} = 3,48 \$
\end{align*}

Cote de l'obligation 

\begin{align*}
cote_k = P_k - IC_k
\end{align*}

\begin{align*}
cote_{0,87} = 85 - 3,48
\end{align*}

\begin{align*}
cote_{0,87} = 81,52
\end{align*}

L'obligation vaut donc 81,52 \% de sa valeur nominale.

\newpage


\section*{Exercice 9}
À l'aide d'une calculatrice financière, trouver le taux de rendement à l’échéance eectif semestriel d’une obligation d’une valeur nominale de 10 000 \$ venant à échéance dans 5 ans et ayant un taux de coupon semestriel de 2,5 \% payable deux fois par année et un prix de 8980 \$.

\begin{enumerate}[a)]
\item 3.4\%
\item 3.5\%
\item 3.6\%
\item \hl{3.7\%}
\end{enumerate}

\section*{Solution 9}
Fixer les paramètres  dans la formule du prix $P_0$ d'une obligation.\\

\textbf{Paramètres:}
\begin{itemize}
\item $P_0 = 8980$ \$
\item $M = 10 000$ \$
\item $C = 0,025 \times 10000 = 250$ \$
\item $n = 10$ semestres
\item $r =$ \?
\end{itemize}

\begin{align*}
P_0=C \times \left[\frac{1-\left( \frac{1}{(1+r)^n}\right)}{r} \right] +\frac{M}{(1+r)^n}
\end{align*}

\begin{align*}
8980=250 \times \left[\frac{1-\left( \frac{1}{(1+r)^{10}}\right)}{r} \right] +\frac{10000}{(1+r)^{10}}
\end{align*}



Avec une calculatrice financière,  entrer les paramètres et estimer $r$ par procédé itératif.  $(r=I/Y)$
\begin{itemize}
\item $10 \rightarrow n$
\item $-8980 \rightarrow PV$
\item$250 \rightarrow PMT$ 
\item $10 000 \rightarrow FV$
\item $COMP \rightarrow I/Y$
\item $ I/Y = 3,74$ \%
\end{itemize}
\textbf{Note:} \\ Il faut mettre le terme PV en négatif dans la calculatrice afin d'obtenir la bonne réponse puisque l'investisseur doit effectuer une sortie de fonds de - 8980 \$ pour acheter l'obligation. Autrement, vous obtiendrez un erreur lors du calcul de I/Y.


\newpage



\section*{Exercice 10}
Trouver le taux de rendement courant d'une obligation ayant un taux de coupon annuel de 10 \% payable deux fois par année, une valeur nominale de 1000 \$ ainsi qu'un prix au marché de $P_0$  = 1100 \$.

\begin{enumerate}[a)]
\item 9,01\%
\item 9,99\%
\item \hl{9,09\%}
\item 9,49\%
\end{enumerate}

\section*{Solution 10}

Taux courant \\

Paramètres:
\begin{itemize}
\item $M = 1000$ \$
\item $P_0 = 1 100$ \$
\item $C_{annuel} = 10 \% \times 1000 \$ = 100 \$$
\end{itemize}

Taux de rendement courant
\begin{align*}
= \frac{C_{annuel}}{P_0}
\end{align*} 

Taux de rendement courant 
\begin{align*}
= \frac{100\$}{1100\$}
\end{align*} 

Taux de rendement courant 
\begin{align*}
= 9,09091 \%
\end{align*} 
\textbf{Note:} \\ Pour le taux de rendement courant,  il faut utiliser les coupons annuels et non pas les coupons semestriels.

\newpage


\section*{Exercice 11}
Un investisseur possède une obligation payant, en deux versements, un coupon annuel de 12 \%.  L’obligation vient à échéance dans 10 ans et promet un taux de rendement annuel nominal de 10 \% (tous les taux nominaux sont capitalisés semestriellement).  L’investisseur a un horizon de 2 ans et prévoit réinvestir ses coupons à un taux annuel nominal de 8 \% pendant cette période.  Quel rendement effectif annuel l’investisseur prévoit-il réaliser s’il croit être en mesure de vendre son obligation à la fin de son horizon à un taux de rendement annuel nominal de 12 \%?

\begin{enumerate}[a)]
\item 0,66\%
\item \hl{0,56\%}
\item 0,60\%
\item 0,50\%
\end{enumerate}

\section*{Solution 11 }
Prix d'achat de l'obligation aujourd'hui:\\

Paramètres: 
\begin{itemize}
\item $r_{achat} = \frac{0.1}{2} = 0.05$
\item $C = \frac{0.12}{2} \times 100 =6$ \$
\item $n = 20$ semestres
\end{itemize}

\begin{align*}
VP=C \times \left[\frac{1-\left( \frac{1}{(1+r)^n}\right)}{r} \right] +\frac{M}{(1+r)^n}
\end{align*}

\begin{align*}
P_0=6 \times \left[\frac{1-\left( \frac{1}{(1+0.05)^{20}}\right)}{0.05} \right] +\frac{100}{(1+0.05)^{20}}
\end{align*}

\begin{align*}
P_0=74.77326+3768895
\end{align*}

\begin{align*}
P_0=112,46221 \$
\end{align*}
\textbf{Note:}\\ Il s'agit du prix que l'investisseur débourse aujourd'hui afin d'acquérir l'obligation.  Autrement dit, c'est la mise de fonds initiale pour effectuer l'investissement


\newpage

Valeur finale des coupons réinvestis sur l'horizon d'investissement: \\


\textbf{Paramètres:}
\begin{itemize}
\item $r_{réinvestissement} = \frac{0.08}{2} = 0.04$
\item $C =6$ \$
\end{itemize}

\begin{align*}
VF=C \times \left[ \frac{(1+r)^n-1}{r} \right]
\end{align*}

\begin{align*}
VF_4= 6 \times \left[ \frac{(1+0.04)^4-1}{0.04} \right]
\end{align*}

\begin{align*}
VF_4=25.47878\$
\end{align*}
Valeur de revente de l'obligation à la fin de l'horizon d'investissement (2 ans):\\
Paramètres: 
\begin{itemize}
\item $r_{revente} = \frac{0.12}{2} = 0.06$
\item $C = \frac{0.12}{2} \times 100=6$\$
\item $n=20-4=16$ semestres 
\item $M=100$ \$
\end{itemize}
\begin{align*}
VP=C \times \left[\frac{1-\left( \frac{1}{(1+r)^n}\right)}{r} \right] +\frac{M}{(1+r)^n}
\end{align*}

\begin{align*}
P_4=6 \times \left[\frac{1-\left( \frac{1}{(1+0.06)^{16}}\right)}{0.06} \right] +\frac{100}{(1+0.06)^{16}}
\end{align*}

\begin{align*}
P_4=60.63537+3936463
\end{align*}

\begin{align*}
P_4=100\$
\end{align*}
\textbf{Note:} \\Puisque l'obligation vient à échéance dans 20 semestres aujourd'hui et que nous la disposons dans 4 semestres, l'obligation continuera de verser des coupons durant 16 semestres (20 - 4) avant d'arriver à terme. Il faut donc actualiser les flux monétaires pour les 16 semestres restants.

\newpage

Calcul du rendement effectif annuel:\\

Paramètres:
\begin{itemize}
\item Prix d"achat de l'obligation $\rightarrow$ $P_0=112.46221$ \$
\item Valeur finale des coupons réinvestis pendant 4 semestres$\rightarrow$ $VF_4=25.47878$\$
\item Prix de vente de l'obligation $\rightarrow$ $P_4=100$ \$
\item $n=2$ ans 
\end{itemize}

\begin{align*}
R= \left[ \frac{\textbf{Produit de disposition de l'obligation}}{\textbf{Mise de fonds initiale}} \right]^{\frac{1}{n}}-1
\end{align*}

\begin{align*}
R= \left[ \frac{VF_4+P_4}{P_0} \right]^{\frac{1}{n}}-1
\end{align*}

\begin{align*}
R= \left[ \frac{25.47878 + 100}{112.46221} \right]^{\frac{1}{2}}-1
\end{align*}

\begin{align*}
R= \left[ \frac{125.47878 }{112.46221} \right]^{\frac{1}{2}}-1
\end{align*}

\begin{align*}
R = 0.056287
\end{align*}
\textbf{Interprétation:} \\ Le rendement effectif annuel pour un horizon d'investissement de 2 ans est d'environ 5,6 \% en utilisant les hypothèses mentionnées dans le problème.
\end{document}